\section{Base Functions}

\label{sec:functions}
% ============================================================================

Base Functions are TypeScript functions that run inside the Base process. They
are not Lambda functions. They are not containers. They are JavaScript closures
executed in an embedded Goja VM, sharing the same process and the same database
connection as the rest of Base.

\subsection{Execution Model}

The Goja VM implements ECMAScript 2023. TypeScript is transpiled to JavaScript
at deploy time (via esbuild, also embedded in the binary). Functions execute
in response to events:

\begin{itemize}[leftmargin=1.1em]
    \item \textbf{Before hooks}: Run before a database operation (create, update,
    delete). Can modify the record or reject the operation.
    \item \textbf{After hooks}: Run after a database operation. Used for
    side effects (notifications, audit logs, cache invalidation).
    \item \textbf{Routes}: Custom HTTP endpoints registered in the Base router.
    Full access to request/response, database, and auth context.
    \item \textbf{Cron}: Scheduled functions executed at configurable intervals.
    Uses the Base internal scheduler, not system cron.
    \item \textbf{Realtime}: Functions triggered by WebSocket subscription
    events. Used for live data transformation.
\end{itemize}

\subsection{Zero Cold Start}

There is no cold start. The Goja VM is initialized once at Base startup and
persists for the lifetime of the process. Function invocation is a direct
Go function call into the VM---no process spawn, no container boot, no
network hop. Measured overhead per invocation: 0.3~ms for a trivial function,
1.2~ms for a function that reads 10 records and writes 1.

\begin{table}[H]
\centering
\small
\begin{tabular}{lrr}
\toprule
\textbf{Operation} & \textbf{Cold Start} & \textbf{Warm} \\
\midrule
AWS Lambda (Node.js) & 200--800~ms & 5--15~ms \\
Cloudflare Workers & 0~ms & 1--3~ms \\
Supabase Edge Functions & 50--200~ms & 3--8~ms \\
Base Functions (Goja) & 0~ms & 0.3--1.2~ms \\
\bottomrule
\end{tabular}
\caption{Function invocation latency comparison. Base Functions have zero cold start because the VM is in-process.}
\label{tab:coldstart}
\end{table}

\subsection{Capability Model}

Functions execute with a capability set, not ambient authority. Each function
declares what it can access:

\begin{itemize}[leftmargin=1.1em]
    \item \textbf{Collections}: Which collections the function can read/write.
    \item \textbf{Network}: Whether the function can make outbound HTTP requests.
    \item \textbf{KMS}: Whether the function can access encrypted fields.
    \item \textbf{Auth}: Whether the function can read/modify user sessions.
\end{itemize}

A function that declares no capabilities can only compute on its inputs and
return a result. This is the default. Capabilities are additive, never
implicit.

% ============================================================================
